\chapter{Method II --- angular identifiability of spherical harmonics}
\label{ch:method2}

The same argument applies on the sphere of directions. Each primitive carries
view-dependent colour as spherical harmonics to degree $3$: $16$ coefficients per
channel, $48$ floats, against $11$ for position, scale, rotation and opacity
combined. Non-DC spherical harmonics are $45$ of $59$ floats --- $76\,\%$ of the
model. They are fitted from however many directions happen to see the primitive,
which may be very few.

Let $y(d)$ be the vector of SH basis functions evaluated at direction $d$. The
angular Gram matrix accumulated over views is
\begin{equation}
  A_k \;=\; \sum_n w_{nk}\, y(d_{nk})\, y(d_{nk})^{\!\top}.
  \label{eq:gram}
\end{equation}
A coefficient block is determinable only where the corresponding block of $A_k$
is well conditioned. The criterion is to retain degree $\ell$ for primitive $k$
only where
\begin{equation}
  \ell_{\max}(k) \;=\; \max\big\{\, \ell \;:\; \lambda_{\min}\big(A_k(\ell)\big) > \tau \,\big\},
  \label{eq:lmax}
\end{equation}
and to mask higher degrees to zero. A primitive seen from three directions
cannot support degree $3$; fitting it produces coefficients that encode noise
and, worse, produce visible artefacts under novel views that fall outside the
observed cone.

\keybox{why this is a separate contribution, not a footnote}{%
B2 is independent of B1. It survives even if the floor-occupancy diagnostic of
\S\ref{sec:instruments} kills the positional band-limit: the two share a
principle --- do not fit what the capture does not determine --- but not a
mechanism, and not a failure mode. B2 also carries a memory result that B1 does
not, since masked coefficients need not be stored.}

\section{Making $\tau$ mean the same thing everywhere}
\label{sec:b2scale}

Equation~\eqref{eq:lmax} as written thresholds $\lambda_{\min}(A_k(\ell))$
against an absolute $\tau$, and that does not survive contact with a real
capture. Two separate scale dependencies are at fault, and they need separate
fixes.

\paragraph{View count.} $A_k$ is a \emph{sum} over views, so a primitive seen by
$100$ cameras has eigenvalues an order of magnitude larger than one seen by $10$,
before any question of conditioning arises. The normalised form
\begin{equation}
  \bar A_k \;=\; \frac{A_k}{\sum_n w_{nk}}
  \label{eq:gramnorm}
\end{equation}
removes it: $\bar A_k$ is a weighted \emph{average} of $y y^{\!\top}$ over the
directions that saw the primitive, so $\tau$ means the same thing for primitives
with different view counts. Without this, $\tau$ silently encodes a preference
for well-observed primitives that has nothing to do with whether their angular
sampling is well spread.

\paragraph{Angular spread.} Normalising is necessary and not sufficient. Even on
$\bar A_k$, $\lambda_{\min}$ still scales with how widely the directions are
spread, so one $\tau$ cannot serve a $100$-camera orbit and a $10$-camera cone.
Section~\ref{sec:b2amend} measures this: at $\tau = 0.01$ the absolute form keeps
everything on a full orbit and nothing on either degraded protocol, and at
$\tau = 0.001$ it keeps $82\,\%$ on the arc. It is a cliff, not a graded
response. The scale-free condition number
\begin{equation}
  \ell_{\max}(k) \;=\; \max\Big\{\, \ell \;:\;
  \frac{\lambda_{\min}\big(\bar A_k(\ell)\big)}{\lambda_{\max}\big(\bar A_k(\ell)\big)} > \tau \,\Big\}
  \label{eq:lmaxcond}
\end{equation}
behaves as the criterion should at a single threshold, and is what the
implementation uses. Equation~\eqref{eq:lmax} is retained above because it is
the form the argument is most naturally stated in, and because the amendment is
a result: it is a change to the method that only measurement could have
produced, and \S\ref{sec:b2amend} reports the measurement that forced it rather
than folding it silently into the derivation.

\paragraph{Status.} B2 is implemented and measured. The criterion is evaluated on
all eight Blender scenes and across the five capture protocols of
\S\ref{sec:stress} without a GPU, since Equation~\eqref{eq:gram} depends only on
primitive positions and training view directions; Table~\ref{tab:b2} adds the
rendered-quality and model-size columns. One caveat carries through to every
statement about memory: masking coefficients to zero does not shrink a
fixed-width PLY, so \emph{no model-size saving is claimed from these numbers}.
Realising it needs variable-length SH storage, which Chapter~\ref{ch:plan}
schedules for Spring.
